\documentclass[11pt,a4paper]{article}

\usepackage[top=25mm,bottom=25mm,left=25mm,right=25mm]{geometry}
\usepackage{kotex}
\usepackage{amsmath,amssymb}
\usepackage{array}
\usepackage{booktabs}
\usepackage{tabularx}
\usepackage{hyperref}
\usepackage{xurl}
\usepackage{enumitem}
\usepackage{setspace}
\usepackage{graphicx}

\setmainhangulfont{AppleMyungjo}
\setsanshangulfont{Apple SD Gothic Neo}
\setlength{\parindent}{0pt}
\setlength{\parskip}{0.45em}

\hypersetup{
  colorlinks=true,
  linkcolor=black,
  citecolor=black,
  urlcolor=blue
}

\setstretch{1.15}
\setlength{\parskip}{0.5em}
\setlength{\parindent}{1.5em}
\emergencystretch=2em

\newcommand{\placeholderfigure}[2]{%
  \fbox{%
    \parbox[c][#1][c]{0.92\linewidth}{%
      \centering
      \textbf{도식 삽입 예정}\par
      \vspace{0.5em}
      \small #2
    }%
  }%
}

\newenvironment{editmemo}[1]
{
\par\noindent\rule{\linewidth}{0.4pt}\par
\noindent\textbf{[편집 메모 시작: #1]}\par
\small\itshape
}
{
\par\normalsize\normalfont
\noindent\textbf{[편집 메모 끝]}\par
\noindent\rule{\linewidth}{0.4pt}\par
}

\title{구간 정보와 긴급차량 접근을 통합한 실시간 주행 경고 컨셉:\\CANoe SIL 기반 구현 및 시험 사례}
\author{현대모비스 PBL 팀}
\date{2026-03-18}

\begin{document}

\maketitle

\begin{abstract}
본 연구는 구간 정보, 긴급차량 접근 정보, 객체 위험 정보가 동시에 유입될 때 경고 우선순위와 복귀 규칙을 일관되게 처리하기 위한 실시간 주행 경고 컨셉을 대상으로 한다. 개별 경고 기능을 분리해 운용할 경우 우선순위 충돌, 표시 중복, timeout 이후 복귀 불안정이 반복되어 이를 하나의 계층형 경고 중재 구조로 통합할 필요가 있었다. 이를 위해 현대 차량 시스템을 참고한 기본 ECU 구성과 핵심 서비스 경로를 CANoe SIL 환경에 재구성하고, CAN과 Ethernet을 공식 통신 범위로 하여 경고 컨셉을 구현하였다. 구현은 CAPL 기반 논리 ECU와 검증 하네스를 중심으로 수행하였으며, V2X 기반 긴급 이벤트 정규화, ADAS 기반 위험도 산정, 경고 중재 및 복귀 규칙을 실행 가능한 로직으로 구성하였다. 또한 요구사항, 기능, 통신, 변수, 코드, 단위시험·통합시험·시스템시험을 연결하는 추적 구조와 HARA, Test Oracle, Test Matrix를 함께 유지하였다. 그 결과 본 연구는 구간 경고, 긴급 경고, 객체 위험 경고를 하나의 중재 구조 안에서 설명하고 반복 시험할 수 있는 기준선을 마련하였으며, V2X 긴급 이벤트 선택 정확도, 교차로·합류구간 통합 경고 성공률, 복귀 안정성 준수율의 세 지표로 정량 결과를 해석할 수 있는 구조를 확보하였다.
\end{abstract}

\noindent\textbf{핵심어}: 실시간 주행 경고, CANoe SIL, CAN/Ethernet, 긴급차량 접근, HMI, Traceability

\section{서론}

\subsection{문제 배경}

최근 차량 경고 기능은 단일 이벤트에 반응하는 표시 수준을 넘어, 주행 맥락, 외부 이벤트, 차량 상태, 표시 채널 특성을 함께 고려하는 통합 형태로 확대되고 있다. 특히 스쿨존, 고속 구간, 유도 구간과 같은 구간 정보와 경찰차 및 구급차 접근 정보가 동시에 존재하는 경우에는 단일 ECU 동작만으로 기능을 설명하기 어렵고, 입력-판단-출력 관계를 시스템 수준에서 정리할 필요가 있다.

본 연구가 이 주제를 선택한 이유는 경고를 발생시키는 개별 원인보다, 동시에 들어온 여러 원인을 어떤 하나의 경고로 정리하고 어떻게 복귀시키는지가 실제 사용자 경험과 안전성에 더 큰 영향을 준다고 판단했기 때문이다. 차량에는 이미 속도, 조향, 제동, 구간 정보, 표시 채널, 서비스 상태와 같은 개별 정보가 존재하지만, 이를 단일 경고 정책으로 중재하는 구조는 별도로 설계하지 않으면 일관성을 유지하기 어렵다. 이에 본 연구는 새로운 센서나 하드웨어를 추가하는 대신, 기존 차량 정보와 외부 입력을 결합하여 하나의 경고 컨셉으로 구조화하는 접근을 택하였다. 프로젝트 개요와 컨셉 설계 문서 역시 같은 문제의식을 바탕으로 ``구간 정보와 긴급차량 접근 정보를 통합하고, TTC 기반 위험 판단과 자동 감속 보조를 결합한다''는 방향을 초기 기준으로 제시한다~\cite{proj_overview,concept_design}.

구현 대상은 특정 단일 ECU 기능이 아니라, 현대 차량 시스템을 벤치마킹한 기본 ECU 구성과 핵심 서비스 경로를 CANoe SIL에서 재현 가능한 형태로 옮기는 것이었다. 따라서 본 연구는 경고 로직의 설계와 시험 구조를 우선 확정하고, 그 안에서 입력-판단-출력 관계가 일관되게 설명되는지를 중심으로 접근하였다.

\subsection{문제 사례 관찰}

본 과제 수행 과정에서는 세 가지 문제가 반복적으로 관찰되었다. 첫째, 구간 경고, 긴급 경고, 객체 기반 위험 경고가 동시에 유효할 수 있어 단순 기능 구현만으로는 우선순위와 복귀 규칙을 설명하기 어려웠다. 둘째, 차량 상태는 주로 CAN을 통해 전달되고 경고 중재와 일부 전달 경계는 Ethernet을 사용하므로, 동일 기능이 혼합 네트워크 경로를 거치면서 설계 설명과 시험 기대값이 쉽게 분리되었다. 셋째, 경고 결과가 앰비언트, 클러스터, 전면 표시, 오디오, 감속 보조 요청으로 동시에 분기되기 때문에, 동일 입력에 대해 시스템 수준의 일관성을 확보하기가 어려웠다.

따라서 본 주제는 새로운 경고를 하나 더 추가하는 문제가 아니라, 여러 입력과 여러 출력 사이에 단일한 해석 규칙을 두는 문제로 정의되었다. 실제로 긴급 경고와 구간 경고가 겹치는 경우에는 단순 표시 여부보다 우선순위, timeout clear, fail-safe, 복귀 안정성이 함께 확인되어야 했고, CAN과 Ethernet이 혼재된 경로에서는 입력 수신과 경고 출력이 모두 존재하더라도 동일한 기능 의미로 해석되는지 별도 검토가 필요했다. 또한 시험 자산이 늘어날수록 단순 통과 여부보다 상위 설계 의도와 시험 기대값이 같은 층을 보고 있는지 점검하는 작업이 중요해졌다.

대표 사례로는 스쿨존 과속 경고 반영 지연, 경찰·구급 경고와 일반 구간 경고의 충돌, timeout clear 이후 복귀 안정성, 전달 경계 이상 시 fail-safe 유지가 있었다. 이러한 사례들은 모두 단일 ECU 수준보다 경고 컨셉의 시스템 경로를 먼저 설명해야 해결 가능한 문제였다.

\subsection{연구 목적}

이에 따라 본 연구는 주행 상황 실시간 경고 컨셉을 차량 시스템 관점에서 재정리하고, CANoe SIL 기반에서 설계와 시험이 함께 설명될 수 있는 구조를 제안하는 데 목적이 있다. 다시 말해, 본 연구의 목적은 차량 전체 시스템을 새로 설계하는 것이 아니라, 차량 시스템에 적용 가능한 경고 컨셉을 선정하고 이를 V2X 기반 긴급 이벤트 정규화, ADAS 기반 위험도 산정, 계층형 경고 중재 구조로 구체화하는 데 있다.

\section{연구 범위와 목표}

본 연구의 범위는 CANoe SIL 환경, CAN + Ethernet 혼합 네트워크, CAPL 기반 논리 ECU 구현, 그리고 문서-코드-시험 간 추적 구조로 한정하였다~\cite{canoe}. 기본 ECU 구성과 핵심 서비스 경로는 현대 차량 시스템을 벤치마킹해 재구성하였으나, 라이선스와 검증 범위 제약에 따라 공식 통신 범위는 CAN과 Ethernet으로 제한하였다. 실차 네트워크, HIL, 물리 ECU, 상용 UDS 스택, 양산 수준의 timing 검증은 본 연구 범위에 포함하지 않았다.

연구 목표는 다음과 같이 정리하였다.

\begin{enumerate}[label=\arabic*.]
\item 구간 정보와 긴급차량 접근 기반 경고를 하나의 경고 컨셉으로 통합한다.
\item 입력, 판단, 출력, 전달 경계, 관측 경로를 분리한 논리 ECU 구조를 정립한다.
\item CAN과 Ethernet이 혼재된 경고 전달 경로를 일관된 기능 경로로 설명 가능하도록 정리한다.
\item 요구사항부터 시험까지 이어지는 추적 구조 위에 실행 가능한 시험 자산을 구축한다.
\item HARA와 추적성 매트릭스를 통해 ISO 26262 및 ASPICE SWE.4/SWE.5 관점의 검증 근거를 연결한다.
\end{enumerate}

여기서 ISO 26262와 ASPICE는 인증 또는 완전 준수를 선언하는 근거가 아니라, 문서 구조와 추적성 구성의 참조 기준으로 사용하였다~\cite{iso26262,aspice_spice,aspice_training}. 프로젝트 내부에서는 \texttt{01}, \texttt{03}, \texttt{0302}, \texttt{0303}, \texttt{0304}, \texttt{04}, \texttt{05}, \texttt{06}, \texttt{07}, \texttt{00d\_HARA\_Worksheet}, \texttt{00g\_Master\_Test\_Matrix}를 주요 기준 문서로 유지하였으며, \texttt{00g\_Master\_Test\_Matrix}는 프로젝트 Test Matrix로 기능하도록 구성하였다.

\section{시스템 설계}

본 장에서는 경고 컨셉의 핵심을 이루는 V2X 긴급 이벤트 정규화, ADAS 위험도 산정, 계층형 경고 중재, 혼합 네트워크 전달 구조를 중심으로 설계 내용을 정리한다.

\subsection{V2X 기반 긴급 이벤트 정규화}

본 컨셉에서 V2X는 긴급 이벤트를 단순 수신하는 ECU가 아니라, 다수의 긴급 접근 입력을 하나의 활성 이벤트로 정규화하는 역할을 담당한다. 경고 판단 이전에 긴급 이벤트를 먼저 정규화한 이유는, 동일 시점에 여러 긴급차량 접근 정보가 존재하더라도 ADAS와 HMI 계층은 최종적으로 하나의 일관된 긴급 맥락만을 참조해야 하기 때문이다. 이 접근은 긴급차량 우선 통행을 다루는 V2X 연구에서 공통적으로 강조되는 ``선행 우선순위 판단 후 하위 서비스로 전달''의 구조와도 부합한다~\cite{v2x_evp_review}. 내부 기능 정의와 개념 설계, 변수 명세 역시 긴급차량 유형, ETA, SourceID 순서로 단일 경고를 결정하는 규칙을 핵심 원칙으로 제시한다~\cite{func_def,concept_design,sysvar_spec}.

시간 $t$에서 수신된 긴급 이벤트 집합을 $\mathcal{E}_t=\{e_1,e_2,\dots,e_n\}$라 하면, 각 이벤트는 차량 유형 $type$, 도착예상시간 $ETA$, 발신 식별자 $SourceID$의 조합으로 표현할 수 있다. 본 연구에서는 긴급차량 유형 우선순위를 먼저 적용하고, 같은 유형 내에서는 ETA가 짧은 후보를 우선하며, ETA가 동일할 경우 SourceID가 작은 후보를 선택하는 사전식 규칙을 사용하였다. 이를 식으로 정리하면 최종 긴급 이벤트는

\[
e_t^*=\arg\min(-rank(type), ETA, SourceID)
\]

로 표현할 수 있다. 여기서 $rank(type)$는 구급차를 경찰차보다 높은 우선순위로 두는 함수이다. 기능 정의서, 시스템 기능 분석서, 시스템 변수 명세는 이러한 우선순위를 ``구급 우선, ETA 우선, SourceID 동률 해소'' 규칙으로 정리하고 있으며~\cite{func_def,sysfunc_analysis,sysvar_spec}, CAPL 구현은 동일 규칙을 기반으로 활성 긴급 이벤트를 유지한다.

여기서 $\arg\min$은 $(-rank(type), ETA, SourceID)$에 대한 사전식 최소 연산을 의미한다. 즉 차량 유형 우선순위를 먼저 적용하고, 동일 유형에서는 ETA가 짧은 후보를, 다시 동일할 경우 SourceID가 작은 후보를 선택한다.

긴급 이벤트는 선택 이후에도 무한정 유지되지 않는다. 활성 이벤트가 일정 시간 내에 갱신되지 않으면 stale로 판단하고, 안전한 해제 상태로 전환한 후 복귀 조건을 다시 평가한다. 본 연구에서는 이 과정을 watchdog 기반 상태 전이로 처리하였다. 이 설계는 통합시험에서 timeout clear와 복귀 시나리오로 직접 다루어지며~\cite{integration_test}, 경고가 단순 발생 여부가 아니라 생성, 유지, 해제까지 포함한 연속 상태임을 보여준다.

\begin{figure}[t]
\centering
\placeholderfigure{48mm}{다중 긴급 이벤트 입력, 차량 유형 우선순위 적용, ETA 비교, SourceID 동률 해소, 활성 이벤트 유지, watchdog timeout clear, 정상 상태 복귀를 포함한 상태도를 배치한다. 좌측은 EmergencyAlert 입력 집합, 중앙은 V2X 정규화, 우측은 활성 긴급 이벤트 $e_t^*$로 구성한다.}
\caption{V2X 긴급 이벤트 선택 및 timeout clear 상태도}
\label{fig:v2x_state}
\end{figure}

\subsection{ADAS 위험도 산정과 계층형 경고 중재}

ADAS는 구간 맥락, 긴급 접근 정보, 객체 기반 위험 정보를 하나의 경고 상태로 통합하는 중심 판단 계층이다. 시간 $t$에서의 입력 상태는 차량 상태 $V_t$, 구간 맥락 $C_t$, 긴급 접근 정보 $E_t$, 객체 위험 정보 $O_t$, 서비스 가용성 $S_t$의 조합으로 두고, 최종 경고 상태를

\[
W_t=f(V_t,C_t,E_t,O_t,S_t)
\]

로 표현하였다. 여기서 $f$는 단일 경고 선택 함수이면서 동시에 감속 보조, fail-safe, stale 입력 처리까지 포함하는 계층형 중재 규칙이다.

긴급 접근 위험도는 ETA, 자차 속도, 접근 방향, 긴급차량 유형을 결합한 가중 합으로 산정하였다. 구현에서는 ETA가 짧을수록, 차량 속도가 높을수록, 전방 접근일수록, 구급차일수록 위험도를 높게 부여하였다. 이를 식으로 쓰면

\[
R_{emg}=0.55\cdot \eta_{score}+speed_{score}+direction_{score}+type_{bias}
\]

와 같이 정리할 수 있다. 이는 프로젝트 개요와 컨셉 설계 문서가 제시한 ``긴급 접근 위험도 + TTC 기반 감속 보조'' 방향과 일치하며~\cite{proj_overview,concept_design}, Euro NCAP AEB Car-to-Car 프로토콜과 NHTSA Crash Avoidance Test Reference Guide가 TTC 기반 경고 및 제동 시점의 평가 항목을 명시적으로 다루는 방향과도 맞닿아 있다~\cite{euroncap_aeb,nhtsa_crash_avoidance}. 또한 차량 통신과 센서 융합 기반 충돌 경고 연구는 이러한 시간-위험도 연계 접근이 일반적이라는 점을 보여준다~\cite{sensors_v2x_collision}. 다만 여기서 사용한 가중치와 임계값은 외부 프로토콜에서 직접 채택한 값이 아니라, 프로젝트 구현 규칙으로 보정한 값으로 해석해야 한다.

구현상 $\eta_{score}$는 ETA가 짧을수록 크게 부여되는 선형 점수이고, $speed_{score}$는 자차 속도 정규화 점수, $direction_{score}$는 접근 방향 점수, $type_{bias}$는 긴급차량 유형에 따른 보정항이다. 따라서 위 식은 하나의 절대 모델이라기보다, 시간적 근접도와 상대 주행 맥락을 함께 반영하기 위한 프로젝트형 위험도 함수로 이해하는 것이 적절하다.

객체 기반 위험도는 TTC(Time-to-Collision)와 객체 신뢰도, 거리, 시나리오 편향을 함께 반영하였다. TTC를 경고 시점의 핵심 지표로 두는 접근은 충돌 경고 분야의 고전적 연구와도 연결된다~\cite{kiefer_ttc}. 여기서 $d$는 객체 거리, $|v_{rel}|$은 절대 상대속도이며, 구현에서는 $d$를 거리 단위 입력, $v_{rel}$을 속도 단위 입력으로 받아 TTC를 ms 단위로 정규화한다. 상대 속도를 $v_{rel}$, 객체 거리를 $d$라 하면 TTC는

\[
TTC=\frac{3600\cdot d}{|v_{rel}|}
\]

로 계산하였다. 이후 TTC 구간별 기본 위험도에 거리 편향, 신뢰도 편향, 교차 충돌 및 끼어들기 편향을 더해 객체 위험도를 구성하였다. 최종 위험도는

\[
R_t=\max(R_{emg},R_{obj})
\]

로 두어, 긴급 접근 위험과 객체 위험 중 더 높은 쪽이 우선 반영되도록 하였다.

감속 보조는 단일 임계값이 아니라 hysteresis 구조로 설계하였다. 위험도가 활성 임계값 $T_{on}$ 이상이면 감속 보조를 요청하고, 해제 임계값 $T_{off}$ 이하로 내려가면 해제하도록 하였다. 즉

\[
u_t=
\begin{cases}
1, & R_t \geq T_{on} \\
0, & R_t \leq T_{off} \\
u_{t-1}, & T_{off}<R_t<T_{on}
\end{cases}
\]

의 형태를 갖는다. 이 구조는 긴급 경고, 객체 경고, 감속 보조가 독립적으로 흔들리는 현상을 줄이고, 동일 입력이 유지되는 동안 동일 경고가 유지되도록 하기 위한 것이다. 컨셉 설계 문서 역시 교차로 및 합류 구간에서 TTC 충돌 위험이 기준 이상이면 자동 감속 보조를 요청하도록 정의하고 있으며~\cite{concept_design}, 기능 정의서와 시스템 기능 분석서는 위험도 임계 초과 시 감속 보조 요청을 생성하고 운전자 개입 시 즉시 해제하도록 정리하고 있다~\cite{func_def,sysfunc_analysis}.

최종 경고 선택은 ``긴급 경고 우선, 비긴급 상태에서는 구간 경고와 객체 경고 비교, 감속 보조 활성 시 경고 강도 상향, stale와 fail-safe 시 최소 안전 경고 보장''의 순서로 수행하였다. 즉 본 연구의 경고 중재는 단순 우선순위 나열이 아니라 위험도 산정, 계층형 선택, 안전 보정이 연속적으로 결합된 구조라고 볼 수 있다. 이 규칙은 기능 정의서의 경고 우선순위 항목~\cite{func_def}, 통합시험의 우선순위 및 timeout clear 항목~\cite{integration_test}, 시스템시험의 연속 시나리오 항목~\cite{system_test}과 같은 판단 기준을 공유한다.

\begin{figure}[t]
\centering
\placeholderfigure{52mm}{차량 상태 $V_t$, 구간 맥락 $C_t$, 긴급 접근 정보 $E_t$, 객체 위험 정보 $O_t$, 서비스 가용성 $S_t$가 위험도 계산 블록으로 수렴하고, $R_{emg}$, $TTC$, $R_{obj}$, $R_t$ 및 hysteresis를 거쳐 최종 경고 상태 $W_t$, 감속 보조, fail-safe 보정으로 분기되는 블록도를 배치한다.}
\caption{ADAS 위험도 산정 및 계층형 경고 중재 블록도}
\label{fig:adas_block}
\end{figure}

\subsection{혼합 네트워크 전달 구조}

네트워크는 CAN 기반 차량 상태 경로와 Ethernet 기반 경고 전달 경로를 함께 사용한다. CAN 경로는 차량 속도, 조향, 제동, 구간 정보와 같은 기초 상태를 수집하는 역할을 맡고, Ethernet 경로는 정규화된 긴급 이벤트와 경고 중재 결과를 전달하는 역할을 맡는다. 따라서 본 컨셉의 통신 구조는 단순히 두 버스를 병행 사용한 것이 아니라, ``상태 수집''과 ``경고 전달''을 기능적으로 분리한 혼합 구조로 해석할 수 있다.

이 구조를 입력 수집 경로, 경고 중재 경로, 출력 분배 경로, 관측 경로로 나누면 전체 흐름을 비교적 명확하게 설명할 수 있다. 입력 수집 경로에서는 차량 상태와 구간 정보가 CAN을 통해 수집되고, 긴급 이벤트는 Ethernet을 통해 유입된다. 경고 중재 경로에서는 정규화된 긴급 이벤트와 위험도 산정 결과가 결합되어 단일 경고 상태가 형성된다. 출력 분배 경로에서는 최종 경고가 앰비언트, 클러스터, 전면 표시, 오디오, 감속 보조 요청으로 분기된다. 관측 경로는 시험과 추적성 확인을 위해 운영 경로와 병행되며, 기능 결과를 직접 바꾸지 않고 상태를 읽는 역할에 머물도록 설계하였다.

이와 같은 구조는 V2X 통신과 차량 내 센서 기반 위험 판단을 함께 사용하여 경고를 구성하는 선행연구와 통합 차량 경고 시스템 사례의 일반 구조와도 부합한다~\cite{sensors_v2x_collision,nhtsa_ivbss}. 동시에 본 프로젝트의 네트워크 플로우와 통신 명세는 경고 정보 전달 정책과 경계 상태를 별도 관리하도록 되어 있어~\cite{comm_spec,nw_flow}, 혼합 네트워크를 단순 메시지 채널이 아니라 경고 컨셉의 기능 경계로 해석할 수 있다.

\begin{figure}[t]
\centering
\placeholderfigure{50mm}{CAN 기반 상태 수집 경로와 Ethernet 기반 긴급 이벤트 유입 경로가 중앙의 V2X/ADAS 블록으로 모이고, 우측으로 BCM, IVI, CLU, HUD, AMP, Decel Assist가 분기되는 혼합 네트워크 전달 구조도를 배치한다. 관측 경로는 운영 경로와 병렬로 별도 표기한다.}
\caption{CAN/Ethernet 혼합 경고 전달 구조도}
\label{fig:network_flow}
\end{figure}

\section{구현 및 검증 구성}

본 장에서는 3장에서 정리한 V2X-ADAS 중심 구조가 CANoe SIL 환경에서 어떻게 실행 자산으로 옮겨졌는지, 그리고 그 자산이 어떤 시험 구조와 연결되었는지를 설명한다.

\subsection{CAPL 기반 구현 구성}

구현은 CANoe SIL 환경에서 CAPL 기반 논리 ECU로 구성하였다~\cite{canoe}. 동력 계층은 EMS, TCU, VCU를 중심으로 차량 기초 거동을 생성하고, 차체 계층은 ESC와 MDPS를 중심으로 제동 및 조향 상태를 제공한다. 그 위에서 V2X는 긴급 이벤트 정규화와 활성 이벤트 유지, ADAS는 위험도 산정과 최종 경고 중재를 담당하며, BCM, IVI, CLU, HUD, AMP는 최종 경고를 시각 및 청각 표면으로 변환한다. SGW, DCM, IBOX, EXT\_DIAG는 서비스 상태와 경계 조건을 해석하고 검증에 필요한 상태 정보를 제공하는 진단·관측 계층으로 구성하였다.

이 구현 구조의 특징은 핵심 경고 로직이 V2X와 ADAS에 집중되어 있다는 점이다. V2X는 긴급 이벤트의 선택과 해제를 담당하고, ADAS는 위험도 계산과 경고 상태 결정을 담당한다. 따라서 본 컨셉의 핵심 구현은 개별 출력 ECU에 있는 것이 아니라, 두 계층 사이의 정보 정규화와 중재 규칙에 있다고 볼 수 있다. 구현 명세 역시 위험 판단, 표시·안내, 동력 상태, 게이트웨이, 검증 하네스를 구분하여 정리하고 있으며~\cite{sw_impl}, 이는 실제 시험 자산과 직접 연결되는 구현 단위로 기능한다.

또한 실행 환경에서는 Ambient, Cluster, Navigation, V2X Ingress, Diagnostic Console 패널 구성을 통해 경고 출력과 상태 관측 결과를 동일한 시험 맥락에서 확인할 수 있도록 하였다. 이는 혼합 네트워크 기반 경고 컨셉이 최종 출력 채널과 관측 채널에서 어떻게 표현되는지를 함께 확인할 수 있게 한다.

\begin{figure}[t]
\centering
\placeholderfigure{48mm}{좌측의 시험 입력 및 시나리오, 중앙의 CANoe SIL CAPL 논리 ECU군과 Validation Harness, 우측의 상태 변수·로그·단위시험/통합시험/시스템시험 결과 판정으로 이어지는 구현 및 시험 구성도를 배치한다.}
\caption{CANoe SIL 구현 및 시험 구성도}
\label{fig:canoe_impl}
\end{figure}

\begin{figure}[t]
\centering
\placeholderfigure{46mm}{Ambient\_TopView, Cluster\_Alert, Navigation\_Alert, V2X\_Ingress, Diagnostic\_Console 대표 화면 또는 병렬 캡처를 배치한다. 좌측에는 경고 출력 채널, 우측에는 V2X 입력 및 Diagnostic Console을 배치하여 입력-출력-관측 관계가 동시에 보이도록 구성한다.}
\caption{대표 CANoe 패널 및 진단 콘솔 구성}
\label{fig:canoe_panels}
\end{figure}

\subsection{추적성 기반 검증 구성}

검증은 단위시험, 통합시험, 시스템시험으로 나누어 구성하였다. 단위시험은 V2X의 긴급 이벤트 유지와 timeout clear, ADAS의 위험도 산정과 경고 선택과 같이 개별 판단 규칙을 다룬다. 통합시험은 우선순위와 ETA/SourceID 규칙, timeout clear와 복귀, 객체 위험과 감속 보조 연계를 다룬다. 시스템시험은 외부 송신 주기, 연속 시나리오, fail-safe 복귀를 포함한 전체 흐름을 다룬다.

추적성은 ``Req $\rightarrow$ Func $\rightarrow$ Flow $\rightarrow$ Comm $\rightarrow$ Var $\rightarrow$ Code $\rightarrow$ 단위시험/통합시험/시스템시험''의 1:1 연결을 원칙으로 삼았다. 이를 위해 요구사항, 기능 정의, 네트워크 흐름, 통신 명세, 시스템 변수, 구현, 단위시험, 통합시험, 시스템시험 문서를 순차적으로 연결하고, 추가로 \texttt{00d\_HARA\_Worksheet}와 \texttt{00g\_Master\_Test\_Matrix}를 통해 상위 위험 항목과 시험 기준을 함께 관리하였다.

특히 본 연구의 시험 구조는 V2X의 우선순위 선택 규칙과 ADAS의 계층형 중재 규칙이 실제 시험 항목으로 직접 이어지도록 구성되었다는 점에서 의미가 있다. 예를 들어 ETA와 SourceID에 따른 긴급 이벤트 선택은 통합시험과 시스템시험에서 반복적으로 검증되며~\cite{integration_test,system_test}, timeout clear와 fail-safe 역시 별도의 연속 시나리오로 유지된다~\cite{integration_test,master_test_matrix,system_test}. 따라서 본 연구의 검증 구성은 단순 기능 나열이 아니라, ``정규화된 긴급 이벤트가 어떤 규칙으로 선택되고, 그 결과가 어떤 경고와 복귀 상태로 이어지는가''를 반복해서 확인하는 구조라고 볼 수 있다.

\begin{table}[t]
\caption{V2X-ADAS 추적 구조 요약}
\label{tab:traceability}
\centering
\footnotesize
\setlength{\tabcolsep}{4pt}
\begin{tabularx}{\linewidth}{>{\raggedright\arraybackslash}p{2.0cm} >{\raggedright\arraybackslash}p{2.5cm} >{\raggedright\arraybackslash}p{3.0cm} >{\raggedright\arraybackslash}p{1.8cm} >{\raggedright\arraybackslash}p{1.5cm} >{\raggedright\arraybackslash}p{1.5cm} >{\raggedright\arraybackslash}X}
\toprule
상위 요구 & 핵심 기능 & 주요 변수 & 주요 CAPL 로직 & 대표 UT & 대표 IT & 대표 ST \\
\midrule
긴급 접근 경고 & 긴급 이벤트 정규화 및 ETA-SourceID 선택 & emergencyType,\newline EtaSeconds,\newline SourceId & V2X & UT\_004 & IT\_006 & ST\_016/017 \\
객체 위험 경고 & TTC 기반 객체 위험 산정 & objectTtcMin,\newline objectRiskClass & ADAS & UT\_006 & IT\_018 & ST\_027--029 \\
감속 보조 및 복귀 & 감속 보조, timeout clear, 복귀 & decelAssistReq,\newline timeoutClear & ADAS-CGW & UT\_005 & IT\_010/011 & ST\_025/026 \\
\bottomrule
\end{tabularx}
\end{table}

\section{결과 및 고찰}

\subsection{구축 결과}

본 연구를 통해 주행 상황 실시간 경고 컨셉의 설계 체계와 실행 가능한 시험 자산을 함께 구축할 수 있었다. 요구사항 문서에서는 기본 경고, 구간 경고, 긴급 경고, 우선순위, fail-safe, 객체 위험, 차량 상태 연계를 포함한 기능 범위를 정의하였다. 기능 정의서에서는 입력, 출력, ECU 동작을 검토자가 바로 이해할 수 있는 형식으로 정리하였고, 네트워크 흐름과 통신 명세에서는 CAN과 Ethernet 혼합 경로를 표준 양식으로 연결하였다. 그 위에 V2X는 긴급 이벤트 선택과 해제, ADAS는 위험도 산정과 최종 경고 중재를 담당하는 중심 판단 계층으로 정리되었다.

구현 단계에서는 CAPL 기반 ECU 로직과 검증 하네스를 통해 반복 실행 가능한 자산을 구성하였다. 특히 V2X는 긴급차량 유형, ETA, SourceID를 기준으로 활성 이벤트를 정규화하고, ADAS는 긴급 접근 위험도, 객체 기반 TTC 위험도, 감속 보조 임계값, stale와 fail-safe 보정을 포함한 단일 경고 상태를 산출하도록 구현하였다. 시험 단계에서는 검증 계약, Test Oracle, Test Matrix, 단위시험·통합시험·시스템시험 자산을 같은 추적 구조 위에 구성하고~\cite{suite_policy}, HARA와 시험 기준을 상위에서 연결하였다. 그 결과 본 프로젝트는 문서, 코드, 시험 자산이 동일한 경고 컨셉을 기준으로 연동되는 시험 기반을 갖추게 되었다.

또한 프로젝트 수행 과정에서 관찰된 문제를 설계 자산과 시험 항목에 직접 반영할 수 있었다. 예를 들어 긴급 이벤트의 ETA/SourceID 동률 해소, timeout clear 이후 복귀, fail-safe 진입 후 최소 경고 유지, 객체 위험 경고와 감속 보조 요청의 동시 판단 등은 모두 시스템 수준 시나리오로 정리되어 시험 항목으로 연결되었다. 이는 구현 편의보다 컨셉의 설명 가능성과 시험 가능성을 우선한 결과라고 볼 수 있다.

본 논문의 결과는 경고 컨셉이 문서·코드·시험 자산으로 연결된 정도와, 그 결과가 어떤 정량 지표로 해석되는지를 함께 보여주는 방식으로 정리하였다. 이에 따라 결과는 V2X 선택 규칙, 교차로·합류구간 통합 경고, timeout clear 및 fail-safe 이후 복귀 안정성의 세 축을 중심으로 제시하였다.

대표적으로 통합시험은 IT\_006에서 우선순위와 ETA/SourceID 규칙을, IT\_009에서 timeout clear와 150ms 이내 복귀를, IT\_018에서 긴급 경고와 TTC 기반 위험의 동시 상황을 다루고 있으며~\cite{integration_test}, 시스템시험은 ST\_016, ST\_017에서 동률 해소 규칙을, ST\_018, ST\_019에서 외부 송신 100ms 주기를, ST\_020, ST\_021에서 무갱신 이후 안전 해제와 이전 유효 경고 복귀를, ST\_045, ST\_046에서 연속 시나리오와 fail-safe 복귀를 각각 다루고 있다~\cite{system_test}. 이처럼 동일한 경고 컨셉이 기능 정의, 통신 경계, 통합시험, 시스템시험에서 반복적으로 대응된다는 점이 본 연구의 구축 성과라고 볼 수 있다.

정량 결과는 세 가지 핵심 지표를 중심으로 정리하였다. 첫째, V2X 긴급 이벤트 선택 정확도 [N1]\%는 IT\_006, ST\_016, ST\_017을 기준으로 유형, ETA, SourceID 규칙이 올바르게 적용되었는지를 보여준다~\cite{integration_test,system_test}. 둘째, 교차로·합류구간 통합 경고 성공률 [N2]\%는 IT\_018, ST\_022, ST\_023을 기준으로 긴급 접근 정보와 TTC 기반 위험이 단일 경고 및 감속 보조 요청으로 수렴하는지를 보여준다~\cite{integration_test,system_test}. 셋째, 복귀 안정성 준수율 [N3]\%은 IT\_009, ST\_035, ST\_036, ST\_045, ST\_046을 기준으로 timeout clear, fail-safe, 정상 복귀가 기준 시간과 상태 일관성을 만족하는지를 보여준다~\cite{integration_test,master_test_matrix,system_test}. 각 지표는 각각 $N1=(S_{sel}/T_{v2x})\times100$, $N2=(S_{int}/T_{int})\times100$, $N3=(S_{rec}/T_{rec})\times100$으로 정의하였다. 여기서 $S_{sel}$은 정상 선택 시나리오 수, $T_{v2x}$는 전체 V2X 선택 시나리오 수, $S_{int}$는 정상 통합 경고 시나리오 수, $T_{int}$는 전체 교차로·합류 시나리오 수, $S_{rec}$는 정상 복귀 시나리오 수, $T_{rec}$는 전체 복귀 시나리오 수를 의미한다.

\begin{table}[t]
\caption{대표 시험 시나리오 요약}
\label{tab:test_scenarios}
\centering
\footnotesize
\setlength{\tabcolsep}{4pt}
\begin{tabularx}{\linewidth}{>{\raggedright\arraybackslash}p{1.5cm} >{\raggedright\arraybackslash}p{1.8cm} >{\raggedright\arraybackslash}p{3.0cm} >{\raggedright\arraybackslash}p{2.4cm} >{\raggedright\arraybackslash}X}
\toprule
시험 레벨 & 시험 ID & 핵심 조건 & 확인 대상 & 기대 결과 \\
\midrule
IT & IT\_006 & 경찰-구급 동시 접근 & ETA, SourceID 규칙 & 단일 긴급 경고 일관 결정 \\
IT & IT\_009 & 1000ms 무갱신 & timeout clear & 안전 해제 후 복귀 \\
IT & IT\_018 & 긴급 접근 + TTC 위험 동시 발생 & 계층형 중재 & 긴급 경고 우선 + 감속 보조 판정 \\
ST & ST\_018, ST\_019 & 외부 송신 주기 & Ethernet 경로 유지 & 100ms 주기 유지 \\
ST & ST\_045, ST\_046 & fail-safe 및 복귀 & 최소 경고 유지와 정상 복귀 & 연속 시나리오 안정성 \\
\bottomrule
\end{tabularx}
\end{table}

\subsection{한계와 해석}

현재 결과는 실행 가능한 구조와 시험 기반을 확보한 단계이지, 모든 실행 동작이 완전히 안정화되었다는 의미는 아니다. 실제 수행 과정에서는 혼합 CAN/Ethernet 전달 경계, owner 분리, observer 경로, 세부 판정값 민감도와 같은 문제가 반복적으로 관찰되었다. 특히 네트워크 전달 경계와 경고 관측 경로가 동시에 존재하는 부분에서는 단순 기능 구현보다 시스템 의미 정리가 더 중요하게 나타났다.

따라서 본 연구의 결과는 완전한 종료 상태보다 설계 가능한 경고 컨셉과 시험 가능한 기준선을 확보한 상태로 해석하는 것이 적절하다. 이는 결과를 과장하기보다, 현재 확보한 설계 자산과 시험 체계, 그리고 남아 있는 기술적 과제를 함께 설명하는 방식에 가깝다.

실제로 HARA는 ``긴급 신호 무갱신 시 경고를 안전하게 해제하고 정상 상태로 복귀해야 한다'', ``경고 정보 전달이 유지되어야 하며, 단절 시 강등 동작이 즉시 적용되어야 한다''고 정리하고 있으며~\cite{hara}, Test Matrix 역시 ST\_045, ST\_046을 통해 연속 주행과 fail-safe 복귀를 별도 항목으로 유지한다~\cite{master_test_matrix}. 이는 남은 과제가 단순 기능 추가보다 경계, 복귀, 전달 정책의 정합성에 더 가깝다는 점을 보여준다.

\section{결론 및 제언}

\subsection{결론}

본 연구는 구간 정보, 긴급차량 접근, 객체 위험을 개별 기능이 아닌 하나의 경고 정책으로 통합하는 주행 상황 실시간 경고 컨셉을 정리하고, 이를 CANoe SIL 기반 혼합 CAN/Ethernet 환경에서 실행 가능한 구조로 구체화하였다. 그 결과 V2X는 긴급 이벤트의 유형, ETA, SourceID를 기준으로 활성 이벤트를 정규화하고, ADAS는 긴급 접근 정보와 TTC 기반 객체 위험을 결합하여 단일 경고와 감속 보조 요청을 결정하는 계층형 구조로 구현될 수 있음을 확인하였다.

또한 요구사항, HARA, 기능 정의, 통신·변수 정의, CAPL 구현, 단위시험·통합시험·시스템시험 자산을 하나의 추적 구조로 연결함으로써, 본 컨셉이 단순 아이디어 수준이 아니라 반복 시험과 구조 해석이 가능한 검증 기준선으로 정리될 수 있음을 확인하였다. 특히 V2X 긴급 이벤트 선택 정확도, 교차로·합류구간 통합 경고 성공률, 복귀 안정성 준수율의 세 지표는 본 컨셉의 정량 결과를 대표하는 핵심 축으로 도출되었다.

\subsection{핵심 결과 요약}

본 연구의 결과는 V2X 정규화, ADAS 위험도 산정, 복귀 안정성 검증의 세 축으로 요약할 수 있다. 각 축은 외부 이벤트의 선택 정확도, 복합 시나리오에서의 단일 경고 수렴, timeout clear 및 fail-safe 이후의 정상 복귀 여부를 통해 정량적으로 해석될 수 있다.

\begin{table}[t]
\caption{핵심 정량 결과 지표}
\label{tab:key_metrics}
\centering
\footnotesize
\setlength{\tabcolsep}{4pt}
\begin{tabularx}{\textwidth}{>{\raggedright\arraybackslash}p{2.9cm} >{\centering\arraybackslash}p{2.8cm} >{\raggedright\arraybackslash}X >{\centering\arraybackslash}p{1.2cm}}
\toprule
지표 & 계산식 & 해석 & 값(\%) \\
\midrule
V2X 긴급 이벤트 선택 정확도 & $N1=(S_{sel}/T_{v2x})\times100$ & 유형-ETA-SourceID 규칙 적용 비율 & [N1] \\
교차로·합류구간 통합 경고 성공률 & $N2=(S_{int}/T_{int})\times100$ & 긴급 접근과 TTC 위험의 단일 경고 수렴 비율 & [N2] \\
복귀 안정성 준수율 & $N3=(S_{rec}/T_{rec})\times100$ & timeout clear, fail-safe, 정상 복귀 준수 비율 & [N3] \\
\bottomrule
\end{tabularx}
\end{table}

\noindent\footnotesize 여기서 $S_{sel}$은 정상 선택 시나리오 수, $T_{v2x}$는 전체 V2X 선택 시나리오 수, $S_{int}$는 정상 통합 경고 시나리오 수, $T_{int}$는 전체 교차로·합류 시나리오 수, $S_{rec}$는 정상 복귀 시나리오 수, $T_{rec}$는 전체 복귀 시나리오 수를 의미한다.\normalsize

\subsection{기대효과}

본 연구의 기대효과는 긴급차량 접근과 객체 위험을 동일한 판단 체계에서 다룰 수 있는 경고 구조를 제시했다는 데 있다. V2X는 외부 이벤트를 유형-ETA-SourceID 기반의 비교 가능한 상태로 정규화하고, ADAS는 이를 TTC 기반 객체 위험과 결합하여 단일 경고와 감속 보조 요청으로 수렴시킨다. 이 구조는 긴급차량 접근, 교차로 측방 위험, 합류구간 충돌위험과 같은 복합 시나리오를 동일한 선택 규칙과 복귀 규칙 위에서 해석하게 하며, 향후 협력인지, 신호교차로 인프라 연계, 군집 이벤트 공유와 같은 확장형 V2X 기능으로 이어질 수 있는 기반을 제공한다.

\subsection{제언}

향후 과제는 다음과 같다.

\begin{enumerate}[label=\arabic*.]
\item V2X, ADAS, CGW 사이의 owner 경계와 observer 경로를 더 정제한다.
\item 혼합 CAN/Ethernet 경로의 실행 파라미터 조정과 TC 안정화를 지속한다.
\item 세부 판정값 민감도를 줄이고 시스템 의미 중심의 시험 기준을 확대한다.
\item 경고-감속 보조 반응 시간과 출력 채널 간 동기 편차를 정량 지표로 사용하기 위해, observer timestamp와 XML export 기반 계측 구조를 추가한다.
\item SIL 기반 evidence를 이후 HIL 또는 실차 연계 검증 단계와 연결할 수 있도록 구조를 보강한다.
\end{enumerate}

\begin{thebibliography}{99}

\bibitem{iso26262}
ISO, ``ISO 26262 road vehicles functional safety,'' available: \url{https://www.iso.org/publication/PUB200262.html} (accessed 2026-03-18).

\bibitem{aspice_spice}
intacs, ``SPICE Center,'' available: \url{https://intacs.info/spice-center} (accessed 2026-03-18).

\bibitem{aspice_training}
intacs, ``Training Center,'' available: \url{https://www.intacs.info/training-center} (accessed 2026-03-18).

\bibitem{canoe}
Vector, ``CANoe --- Comprehensive Development, Test and Analysis Environment,'' available: \url{https://www.vector.com/int/en/products/products-a-z/software/canoe/} (accessed 2026-03-18).

\bibitem{func_def}
현대모비스 PBL 팀, ``03\_Function\_definition.md,'' 버전 4.39, 2026-03-17.

\bibitem{integration_test}
현대모비스 PBL 팀, ``06\_Integration\_Test.md,'' 버전 4.23, 2026-03-17.

\bibitem{comm_spec}
현대모비스 PBL 팀, ``0303\_Communication\_Specification.md,'' 버전 3.36, 2026-03-17.

\bibitem{hara}
현대모비스 PBL 팀, ``00d\_HARA\_Worksheet.md,'' 버전 1.6, 2026-03-11.

\bibitem{master_test_matrix}
현대모비스 PBL 팀, ``Test Matrix (00g\_Master\_Test\_Matrix.md),'' 버전 0.2, 2026-03-15.

\bibitem{system_test}
현대모비스 PBL 팀, ``07\_System\_Test.md,'' 버전 5.23, 2026-03-17.

\bibitem{sw_impl}
현대모비스 PBL 팀, ``04\_SW\_Implementation.md,'' 버전 2.26, 2026-03-17.

\bibitem{unit_test}
현대모비스 PBL 팀, ``05\_Unit\_Test.md,'' 버전 2.26, 2026-03-17.

\bibitem{nw_flow}
현대모비스 PBL 팀, ``0302\_NWflowDef.md,'' 버전 3.31, 2026-03-17.

\bibitem{v2x_evp_review}
Ahmed Kamel et al., ``Vehicle-to-Everything (V2X) Communication for Emergency Vehicle Priority: A Comprehensive Review,'' available: \url{https://www.kamelrobotics.com/assets/projects/v2x-evp-review/Vehicle_to_Everything__V2X__Communication_for_Emergency_Vehicle_Priority__A_Comprehensive_Review.pdf} (accessed 2026-03-19).

\bibitem{euroncap_aeb}
Euro NCAP, ``TEST PROTOCOL --- AEB Car-to-Car systems v4.3.1,'' available: \url{https://www.euroncap.com/media/80155/euro-ncap-aeb-c2c-test-protocol-v431.pdf} (accessed 2026-03-19).

\bibitem{nhtsa_crash_avoidance}
National Highway Traffic Safety Administration, ``Crash Avoidance Test Reference Guide, Volume IV, Version 1: Crash Avoidance Database,'' available: \url{https://www.nhtsa.gov/sites/nhtsa.gov/files/2025-05/crash-avoidance-test-reference-guide-volume-4-2024.pdf} (accessed 2026-03-19).

\bibitem{sensors_v2x_collision}
Sensors, ``Vehicle Trajectory Prediction and Collision Warning via Fusion of Multisensors and Wireless Vehicular Communications,'' available: \url{https://www.mdpi.com/1424-8220/20/1/288} (accessed 2026-03-19).

\bibitem{nhtsa_ivbss}
National Highway Traffic Safety Administration, ``Integrated Vehicle-Based Safety Systems Light Vehicle Field Operational Test Final Program Report,'' available: \url{https://rosap.ntl.bts.gov/view/dot/3973/dot_3973_DS1.pdf} (accessed 2026-03-19).

\bibitem{proj_overview}
현대모비스 PBL 팀, ``00\_Project\_Overview.md,'' 버전 0.5, 2026-03-17.

\bibitem{concept_design}
현대모비스 PBL 팀, ``02\_Concept\_design.md,'' 버전 0.5, 2026-03-17.

\bibitem{sysfunc_analysis}
현대모비스 PBL 팀, ``0301\_SysFuncAnalysis.md,'' 버전 3.16, 2026-03-17.

\bibitem{sysvar_spec}
현대모비스 PBL 팀, ``0304\_System\_Variables.md,'' 버전 3.10, 2026-03-17.

\bibitem{kiefer_ttc}
R. J. Kiefer et al., ``Developing an inverse time-to-collision crash alert timing approach,'' \emph{Accident Analysis and Prevention}, vol. 37, no. 2, pp. 295--303, 2005. Available: \url{https://doi.org/10.1016/j.aap.2004.09.003} (accessed 2026-03-19).

\bibitem{suite_policy}
현대모비스 PBL 팀, ``canoe/tests/modules/test\_suites/README.md,'' active suite policy snapshot, 2026-03-21.

\end{thebibliography}

\end{document}
